% Section 2: Method
% Following Arpit's method structure (from editorial principles)

\section{Method}
\label{sec:method}

% Overview paragraph (elevator pitch)
We study what transfers when migrating key-value cache state across large language models.
Our approach decomposes KV injection into four arms---Self, K-only, V-only, and Joint---to isolate the contribution of each component.
We measure task performance via answer log-likelihood and exact match,
and characterize transferability via CCA correlation and layer-wise analysis.

% 2.1 Problem Setup
\subsection{Problem Setup}

Consider a teacher model $T$ and a student model $S$.
Given a document $d$, the teacher produces KV states $(K_T, V_T)$
and the student produces $(K_S, V_S)$.
We ask: which components of $(K_T, V_T)$ can replace $(K_S, V_S)$
while preserving or improving task performance?

\textbf{Key insight.} Keys encode attention routing (addressing);
values encode content representations (semantics).
These may transfer differently across models.

% 2.2 Task and Data
\subsection{Task and Data}

\textbf{Synthetic OOD task.}
We construct a synthetic out-of-domain (OOD) task with novel entities
that do not appear in model pretraining data.
The task requires reasoning across 1--4 hops over a knowledge graph,
with answers verifiable by exact match.

\textbf{Data splits.}
Training set: 42 samples; evaluation set: 14 samples.
We split data by entity clusters to prevent information leakage:
entities in the evaluation set do not appear in the training set.

\textbf{Why synthetic?}
Real-world tasks require ground-truth reasoning chains.
Our synthetic task provides verifiable answers while controlling
entity familiarity---critical for studying OOD transfer.

% 2.3 Model Pairs
\subsection{Model Pairs}

We study 6 pairs from the Qwen3 family:
(8B$\to$0.6B), (4B$\to$1.7B), (4B$\to$0.6B),
(8B$\to$1.7B), (1.7B$\to$0.6B), (8B$\to$4B).

This design covers three configurations:
\begin{itemize}
  \item \textbf{Equal-layer}: 28$\to$28 (4B$\to$1.7B, 1.7B$\to$0.6B)
  \item \textbf{Unequal-layer}: 36$\to$28 (8B$\to$4B), 36$\to$28 (8B$\to$0.6B)
  \item \textbf{Size ratio}: 2.4$\times$ (4B$\to$1.7B) to 13.3$\times$ (8B$\to$0.6B)
\end{itemize}

% 2.4 Injection Protocol
\subsection{Injection Protocol}

For each evaluation sample, we construct four KV configurations:

\textbf{Self.}
Student prefills document with its own KV states.
This is the baseline (doc-only KV).

\textbf{K-only.}
Teacher K (mapped) + Student V.
Tests whether addressing geometry transfers.

\textbf{V-only.}
Student K + Teacher V (mapped).
Tests whether content representations transfer.

\textbf{Joint.}
Teacher K (mapped) + Teacher V (mapped).
Tests full KV transfer (both components).

% 2.5 Mapper
\subsection{Mapper}

We use Affine mappers fitted on calibration data.
The mapper learns a linear transformation $W$ and bias $b$:
\[
K_S = W \cdot K_T + b
\]
For K, we apply de-RoPE before mapping and re-RoPE after.
For V, no positional alignment is needed.

We also tested Whitened and RidgePerHead mappers;
results are invariant to mapper choice (§5).

% 2.6 Metrics
\subsection{Metrics}

\textbf{Primary: Answer log-likelihood.}
We compute teacher-forced log-likelihood of answer tokens.
This is continuous, sensitive to small changes,
and correlates with task performance.

\textbf{Secondary: Exact match.}
We compute greedy exact match on final answers.
This is binary, interpretable, but less sensitive.

\textbf{Baselines.}
\begin{itemize}
  \item \textbf{Self}: Doc-only KV (negative baseline).
  \item \textbf{Student full}: Doc + query KV (upper bound for single model).
  \item \textbf{Teacher full}: Doc + query KV (reference for teacher capability).
\end{itemize}

% 2.7 Statistical Analysis
\subsection{Statistical Analysis}

We report 95\% bootstrap confidence intervals (1000 resamples)
and paired Wilcoxon signed-rank tests ($n = 14$ evaluation samples).
We consider $p < 0.05$ statistically significant.

\textbf{Warning.} With $n = 14$, statistical power is limited.
Wide confidence intervals should be interpreted cautiously.
Larger evaluation sets are needed for definitive conclusions.
